\documentclass[troisieme]{fiche}
\metadonnees{6}{S'envoler}{Bloc B — Ce qu'on peut, ce qu'on doit}

\begin{document}
\entete

\objectifs{%
\textbf{Langue} : \emph{will} et \emph{be going to} ; questions et réponses courtes au
futur.\par
\textbf{Texte} : J.\,M. Barrie, \emph{Peter Pan} (1911), ch. \textsc{iii}.\par
\textbf{Durée} : 1 h — exercices 20 min, texte et questions 25 min, version 15 min.}

%% ===================================================================
\exercice[12 min]{\emph{Will} ou \emph{be going to}}

\rappel{\textbf{Rappel.} \emph{Will} sert à décider sur le moment, à promettre ou à prédire.
\emph{Be going to} annonce une intention déjà prise ou ce qu'un indice rend certain.}

\consigne{Complétez avec \emph{will} ou la forme de \emph{be going to} qui convient.}

\begin{anglais}
\begin{enumerate}[label=\arabic*.,itemsep=2.5mm,leftmargin=*]
  \item ``I can't fly.'' --- ``I \trou{will} teach you.'' \emph{(promesse, sur le moment)}
  \item Look at those clouds: it \trou{is going to} rain.
  \item We \trou{are going to} leave tonight; everything is ready.
  \item ``The window is open.'' --- ``I \trou{will} close it.''
  \item Wendy \trou{is going to} tell the boys a story; she has promised.
  \item One day this boy \trou{will} be famous.
  \item ``Where are you going?'' --- ``I \trou{am going to} tell the other boys.''
  \item Don't worry, I \trou{will} help you with your bag.
  \item She is very tired: she \trou{is going to} fall asleep.
  \item I think they \trou{will} like the story.
\end{enumerate}
\end{anglais}

\note{\emph{Will} : 1, 4, 6, 8, 10. \emph{Be going to} : 2, 3, 5, 7, 9.}

\note[Comment choisir]{Si la décision est prise en parlant (1, 4, 8) ou s'il s'agit d'une
opinion sur l'avenir (6, 10), c'est \emph{will}. Si la décision était prise avant (3, 5, 7)
ou si quelque chose se voit déjà (2, 9), c'est \emph{be going to}. Dans le doute, demandez-vous
s'il y a un indice sous les yeux.}

%% ===================================================================
\exercice[8 min]{Questions et réponses courtes}

\rappel{\textbf{Rappel.} À la question, l'auxiliaire passe devant le sujet. La réponse courte
reprend l'auxiliaire seul, jamais le verbe.}

\consigne{Complétez.}

\begin{anglais}
\begin{enumerate}[label=\arabic*.,itemsep=2.5mm,leftmargin=*]
  \item \trou{Will} you teach me to fly? --- Yes, \trou{I will}.
  \item \trou{Are} they going to leave tonight? --- No, \trou{they aren't}.
  \item Will Wendy come with them? --- No, \trou{she won't}.
  \item \trou{Is} he going to tell the story? --- Yes, \trou{he is}.
  \item Are you going to be afraid? --- No, \trou{I'm not}.
  \item Will the boys listen to her? --- Yes, \trou{they will}.
\end{enumerate}
\end{anglais}

\note[Deux formes à connaître]{La négation de \emph{will} est \emph{will not}, contractée en
\emph{won't} — forme irrégulière, à retenir telle quelle. Et la réponse courte ne répète
jamais le verbe : \emph{Yes, I will}, jamais \emph{Yes, I will teach}.}

%% ===================================================================
\newpage
\section{Texte}

\noindent\small\textit{Peter Pan s'est introduit la nuit dans la chambre des enfants Darling.
Wendy, réveillée, lui a parlé ; il vient de lui dire qu'il ne connaît aucune histoire.}\normalsize

\begin{texte}
``Do you know,'' Peter asked ``why \gl[a swallow]{swallows}{une hirondelle} build in the
\gl[the eaves]{eaves}{l'avant-toit} of houses? It is to listen to the stories. O Wendy, your
mother was telling you such a lovely story.''

``Which story was it?''

``About the prince who couldn't find the lady who wore the glass
\gl[a slipper]{slipper}{une pantoufle}.''

``Peter,'' said Wendy excitedly, ``that was Cinderella, and he found her, and they lived
happily ever after.''

Peter was so glad that he rose from the floor, where they had been sitting, and hurried to the
window.

``Where are you going?'' she cried with \gl[a misgiving]{misgiving}{une appréhension}.

``To tell the other boys.''

``Don't go Peter,'' she \gl[to entreat]{entreated}{supplier}, ``I know such lots of stories.''

Those were her precise words, so there can be no denying that it was she who first tempted
him.

He came back, and there was a \gl[greedy]{greedy}{avide} look in his eyes now which ought to
have alarmed her, but did not.

``Oh, the stories I could tell to the boys!'' she cried, and then Peter
\gl[to grip]{gripped}{empoigner} her and began to draw her toward the window.

``Let me go!'' she ordered him.

``Wendy, do come with me and tell the other boys.''

Of course she was very pleased to be asked, but she said, ``Oh dear, I can't. Think of mummy!
Besides, I can't fly.''

``I'll teach you.''

``Oh, how lovely to fly.''

``I'll teach you how to jump on the wind's back, and then away we go.''

``Oo!'' she exclaimed \gl[rapturously]{rapturously}{avec ravissement}.

``Wendy, Wendy, when you are sleeping in your silly bed you might be flying about with me
saying funny things to the stars.''

``Oo!''

``And, Wendy, there are \gl[a mermaid]{mermaids}{une sirène}.''

``Mermaids! With tails?''

``Such long tails.''

``Oh,'' cried Wendy, ``to see a mermaid!''

He had become frightfully \gl[cunning]{cunning}{rusé}. ``Wendy,'' he said, ``how we should all
respect you.''

She was \gl[to wriggle]{wriggling}{se tortiller} her body in distress. It was quite as if she
were trying to remain on the nursery floor.

But he had no pity for her.

``Wendy,'' he said, the \gl[sly]{sly}{malin, sournois} one, ``you could
\gl[to tuck in]{tuck}{border (dans un lit)} us in at night.''

``Oo!'' and her arms went out to him.

``And you could \gl[to darn]{darn}{repriser} our clothes, and make pockets for us. None of us
has any pockets.''

How could she resist. ``Of course it's awfully fascinating!'' she cried. ``Peter, would you
teach John and Michael to fly too?''

``If you like,'' he said indifferently.
\end{texte}
\source{J.\,M. Barrie, \emph{Peter Pan and Wendy}, Londres, Hodder \& Stoughton, 1911,
ch.~\textsc{iii}.}

\partiel[15 min]{Questions}
\consigne{Répondez aux deux premières questions en français, aux deux dernières en anglais.}

\begin{enumerate}[label=\arabic*.,itemsep=2.5mm,leftmargin=*]
  \item Pourquoi Peter se précipite-t-il vers la fenêtre, et pourquoi revient-il ?
    \lignes{2}
    \corr{Il veut aller raconter aux autres garçons l'histoire de Cendrillon, dont Wendy
    vient de lui donner la fin. Il revient parce que Wendy le supplie de rester et lui dit
    qu'elle connaît des quantités d'histoires.}
  \item Quels arguments Peter emploie-t-il, l'un après l'autre, pour décider Wendy ?
    \lignes{3}
    \corr{Voler avec lui la nuit, les sirènes à longue queue, le respect que tous lui
    porteraient, le fait qu'elle pourrait les border le soir et repriser leurs vêtements. Il
    passe du rêve à la tâche domestique, et c'est le dernier argument qui la décide.}
  \item \en{The narrator says Wendy tempted Peter first. Do you agree? Explain.}
    \lignes{3}
    \corr{She does say \emph{Don't go Peter, I know such lots of stories}, so she offers
    something. But the narrator is being unfair: Peter grips her, pulls her towards the
    window and uses every argument he can find. The word \emph{tempted} hides who is really
    doing the work.}
  \item \en{Find in the text two sentences where Peter promises something with \emph{will},
    and one where Wendy says what she cannot do.}
    \lignes{2}
    \corr{\emph{I'll teach you} ; \emph{I'll teach you how to jump on the wind's back}. Et
    \emph{I can't. Think of mummy! Besides, I can't fly.}}
\end{enumerate}

\note[Le mot qui trahit]{\emph{There was a greedy look in his eyes now which ought to have
alarmed her, but did not.} Le narrateur sait déjà ce que Wendy ne sait pas encore. Toute la
suite du livre tient dans ce \emph{but did not}.}

%% ===================================================================
\section{Version}
\consigne{Traduisez le passage en français. Il précède le texte de quelques lignes : Peter
vient de retrouver son ombre et il pleure.}

\begin{version}
His \gl[a sob]{sobs}{un sanglot} woke Wendy, and she sat up in bed. She was not alarmed to see
a stranger crying on the \gl[a nursery]{nursery}{une chambre d'enfants} floor; she was only
pleasantly interested.

``Boy,'' she said \gl[courteously]{courteously}{courtoisement}, ``why are you crying?''

Peter could be exceeding polite also, having learned the grand manner at fairy ceremonies, and
he rose and \gl[to bow]{bowed}{s'incliner} to her beautifully. She was much pleased, and bowed
beautifully to him from the bed.

``What's your name?'' he asked.

``Wendy Moira Angela Darling,'' she replied with some satisfaction. ``What is your name?''

``Peter Pan.''
\end{version}
\source{\emph{Ibid.}}

\lignes{12}

\corr{\textbf{Proposition de traduction.} Ses sanglots réveillèrent Wendy, qui se redressa
dans son lit. Elle ne s'alarma pas de voir un inconnu pleurer sur le plancher de la chambre
d'enfants ; elle en fut seulement agréablement intriguée.

\og Garçon, dit-elle courtoisement, pourquoi pleures-tu ? \fg

Peter savait être extrêmement poli, lui aussi, ayant appris les belles manières aux
cérémonies des fées : il se leva et s'inclina devant elle magnifiquement. Elle en fut très
flattée, et s'inclina magnifiquement devant lui du fond de son lit.

\og Comment t'appelles-tu ? \fg{} demanda-t-il.

\og Wendy Moira Angela Darling, répondit-elle avec une certaine satisfaction. Et toi, comment
t'appelles-tu ? \fg

\og Peter Pan. \fg}

\note[Choix de traduction 1]{\emph{What's your name?} ne se traduit pas mot à mot :
\og Comment t'appelles-tu ? \fg{} et non \og Quel est ton nom ? \fg, qui sonne comme un
interrogatoire.}

\note[Choix de traduction 2]{\emph{she was only pleasantly interested} : l'adverbe porte sur
l'adjectif. Le français préfère un nom ou un verbe (\og elle en fut agréablement
intriguée \fg) au calque \og seulement agréablement intéressée \fg.}

\note[Choix de traduction 3]{\emph{having learned the grand manner} : participe présent
composé, valeur de cause. En français, \og ayant appris \fg{} ou \og parce qu'il avait
appris \fg. Les deux marchent ; le participe est plus léger.}

%% ===================================================================
\partiel{Lexique à mémoriser}
\begin{multicols}{2}\small
\begin{itemize}[label=--,itemsep=0.8mm,leftmargin=*]
  \item \textbf{a sob} — un sanglot
  \item \textbf{to wake} — réveiller ; woke, woken
  \item \textbf{a stranger} — un inconnu
  \item \textbf{to bow} — s'incliner
  \item \textbf{lots of} — des quantités de
  \item \textbf{to hurry} — se dépêcher \textit{(fiche 1)}
  \item \textbf{greedy} — avide, gourmand
  \item \textbf{to draw} — tirer ; drew, drawn
  \item \textbf{besides} — en outre, d'ailleurs
  \item \textbf{a tail} — une queue
  \item \textbf{cunning} — rusé
  \item \textbf{a pocket} — une poche
\end{itemize}
\end{multicols}

\end{document}
